\documentclass[10pt,a4paper]{amsart}
\usepackage[margin=19mm]{geometry}
\usepackage{amsmath,amssymb,mathtools}
\usepackage[scr=rsfs,cal=euler]{mathalfa}
\usepackage{xfrac}
\usepackage[unicode,hidelinks,bookmarks=false]{hyperref}
\newcommand{\ie}{i.e.\ }
\newcommand{\cf}{cf.\ }
\newcommand{\resp}{respectively\ }
\newcommand{\Subsection}{\subsection}
\providecommand{\nobreakdash}{\nobreak\hbox{-}}
\setlength{\emergencystretch}{2em}
\multlinegap0pt
\usepackage{bm}
\usepackage{bbm}
\usepackage{dsfont}
\usepackage{xspace}
\usepackage{cancel}
\usepackage{mathtools}
\usepackage{amsthm}
\makeatletter
\@ifundefined{thm}{\newtheorem{thm}{Thm}}{}
\@ifundefined{lem}{\newtheorem{lem}[thm]{Lemma}}{}
\@ifundefined{prop}{\newtheorem{prop}[thm]{Proposition}}{}
\makeatother
\makeatletter
\newcommand{\psh}{\operatorname{PSH}}
\expandafter\ifx\csname sciabstract\endcsname\relax
\newenvironment{sciabstract}{}

%%%%%%%%%%%%%%%%% END OF PREAMBLE %%%%%%%%%%%%%%%%


\renewcommand*\contentsname{}
\fi
\makeatother
\title[Metric entropy of the space of holomorphic functions]{Metric entropy of the space of holomorphic functions}
\author{Siarhei Finski}
\date{Source: arXiv:2607.18890v1 (CC BY 4.0)}
\begin{document}
\maketitle

\noindent\emph{Source and licence.} Metric entropy of the space of holomorphic functions. Version arXiv:2607.18890v1 (CC BY 4.0), distributed under the Creative Commons Attribution 4.0 International licence, as recorded in the arXiv licence metadata for this exact version and shown on the abstract page https://arxiv.org/abs/2607.18890v1. Full rights notice: licenses/arxiv-2607.18890-CC-BY-4.0.txt.\par\smallskip

\noindent\textit{Editorial scope.} Counted unit: the Prop whose source label is \texttt{None} in the article (source0.tex lines 1229--1233, source label \texttt{None}), delivered together with the article's own complete original proof of it (source0.tex lines 1234--1272, one \texttt{proof} environment of 39 lines). No mathematics is written, repaired or re-derived: statement and proof are the article's own text line for line. Required context retained uncounted: lem\_extr\_fun\_lem2 (lines 1110--1112); eq\_part\_def\_u\_psh (lines 1127--1137); eq\_vck\_low\_bnd (lines 1155--1163); eq\_rel\_extr\_f (lines 1181--1183); eq\_u\_cnt\_ur (lines 1197--1207). Every definition, display and label of those lines is retained unchanged. Omitted from this package and not redistributed: 1--1109, 1113--1126, 1138--1154, 1164--1180, 1184--1196, 1208--1228, 1273--1487. Nothing inside a retained excerpt is omitted or altered: every retained body is delivered exactly as the article's lines are written, so no omission receipt of any kind accompanies this package. Citations: the retained excerpts carry no citation command at all, so no bibliography entry of the article is transcribed and no reference is redistributed. Reference adaptation: none was needed and none was made; every reference inside the delivered unit resolves inside it, so no unresolved reference or citation is produced. Printed slips kept literal: no printed word, symbol, hypothesis or number of the counted statement or of its proof was altered. Layout and class: the article's own class is \texttt{article} with packages including times, dsfont, enumitem, inputenc, commath, amsthm, amscd, amsmath, amssymb, cite, setspace, stmaryrd, tikz-cd, tocloft; the wrapper is the lane's pinned \texttt{amsart} editorial wrapper at 10pt on A4, which restates the theorem-like environments the retained text needs and adds the article's own macro definitions that the retained text needs (\texttt{\textbackslash psh}), with extra wrapper packages (bm, bbm, dsfont, xspace, cancel, mathtools). The delivered numbering is the unit's own. Assets: no retained span contains a figure environment, a graphics-inclusion command, a PDF-page inclusion, a table or a TikZ picture; no image file is copied into this package, the package declares no asset dependency and no image file is redistributed with it. Licence: version arXiv:2607.18890v1 of this article is distributed under the Creative Commons Attribution 4.0 International licence, as recorded in the arXiv licence metadata for this exact version and shown on the abstract page https://arxiv.org/abs/2607.18890v1. A full rights notice is delivered at licenses/arxiv-2607.18890-CC-BY-4.0.txt. Characters: every retained excerpt is pure ASCII, so the delivered engine, XeLaTeX with its default Latin Modern text font, reports no missing character.
	\begin{lem}\label{lem_extr_fun_lem2}
		For any $0 < r < 1$, there is $v \in \mathscr{L}(X)$ such that $v \leq 0$ in a neighborhood of $\overline{\Omega}$, but $v \geq \delta$ for some $\delta > 0$ in a neighborhood of $\partial \Omega_r$.
	\end{lem}

		\begin{equation}\label{eq_part_def_u_psh}
			u(x) 
			:=
			\begin{cases}
				c \rho(x), \quad &\text{if } x \in \Omega,
				\\
				\max \big\{ v - \epsilon,  c \rho(x) \big\}, \quad &\text{if } x \in \Omega_{1/2} \setminus \Omega,
				\\
				v - \epsilon, \quad &\text{if } x \in X \setminus \Omega_{1/2}.
			\end{cases} 
		\end{equation}

		\begin{equation}\label{eq_vck_low_bnd}
			V_{c K, \Omega}
			\geq
			\begin{cases}
				c \rho(x), \quad &\text{if } x \in \Omega,
				\\
				v - \epsilon, \quad &\text{otherwise.}
			\end{cases} 
		\end{equation}

	\begin{lem}\label{eq_rel_extr_f}
		As $r \to 0$ from above, $u_{K, \Omega_r}$ converges to $u_{K, \Omega}$ pointwise in $\Omega$.
	\end{lem}

		\begin{equation}\label{eq_u_cnt_ur}
			u(x) 
			:=
			\begin{cases}
				v - \epsilon, \quad &\text{if } x \in \Omega_{-\delta},
				\\
				\max \big\{ A (\rho - r),  v - \epsilon \big\}, \quad &\text{if } x \in \Omega \setminus \Omega_{-\delta},
				\\
				A(\rho - r), \quad &\text{if } x \in \Omega_r \setminus \Omega.
			\end{cases} 
		\end{equation}

	\begin{prop}
		For any $x \in \Omega$, the function $c \mapsto u_c(x)$ is decreasing.
		There is $C > 0$ such that $u_c \geq C \cdot \rho$ for any $c > 0$ small enough.
		Finally, as $c \to 0$ from above, $u_c$ converges to $u_{K, \Omega}$ pointwise on $\Omega$.
	\end{prop}

	\begin{proof}
		Note that for any $0 < \alpha < 1$, $\alpha V_{cK, \Omega}$ is one of the competitors in the definition of $V_{\alpha cK, \Omega}$, and so we have the inequality $\alpha V_{cK, \Omega} \leq V_{\alpha cK, \Omega}$, which implies the first claim.
		The second claim is an immediate consequence of (\ref{eq_vck_low_bnd}).
		\par 
		Let us now establish the last claim.
		For this, by Lemma \ref{eq_rel_extr_f}, it suffices to show that for any $r > 0$, there is $c_0 > 0$ such that for any $0 < c < c_0$, we have 
		\begin{equation}\label{eq_vck_cukr}
			V_{cK, \Omega} \geq c u_{K, \Omega_r}.
		\end{equation}
		By Lemma \ref{lem_extr_fun_lem2}, consider $\delta > 0$ and $v \in \mathscr{L}(X)$ so that $v \leq 0$ in a neighborhood of $\overline{\Omega}$, but $v \geq \delta$ in a neighborhood of $\partial \Omega_{r/2}$.
		Let $v_0 \in \psh(\Omega_r)$ be such that $-1 \leq v_0 \leq 0$ and $v_0|_K \leq -1$.
		For any $c > 0$, we define
		\begin{equation}
			u(x) 
			:=
			\begin{cases}
				c v_0, \quad &\text{if } x \in \Omega,
				\\
				\max \big\{ v - \delta/2,  c v_0 \big\}, \quad &\text{if } x \in \Omega_{r/2} \setminus \Omega,
				\\
				v - \delta/2, \quad &\text{if } x \in X \setminus \Omega_{r/2}.
			\end{cases} 
		\end{equation}
		We claim that for $c > 0$ small enough (independent of $v_0$), $u \in \mathscr{L}(X)$. 
		As after (\ref{eq_part_def_u_psh}), it suffices to show that 
		\begin{equation}
			u(x) 
			=
			\begin{cases}
				c v_0, \quad &\text{for $x$ in a neighborhood of $\partial \Omega$},
				\\
				v - \delta/2, \quad &\text{for $x$ in a neighborhood of $\partial \Omega_{r/2}$.}
			\end{cases} 
		\end{equation}
		Note that the first condition is always verified as long as $c < \delta/2$.
		The second condition is also verified, since $v_0 \leq 0$ and $v \geq \delta$ in a neighborhood of $\partial \Omega_{r/2}$.
		As $u$ is one of the competitors for $V_{cK, \Omega}$, we deduce by the maximality that $V_{cK, \Omega} \geq c v_0$ on $\Omega$.
		By an argument similar to the one after (\ref{eq_u_cnt_ur}), we deduce (\ref{eq_vck_cukr}).
	\end{proof}

\end{document}
